% Part: proof-theory
% Chapter: sequent-calculus
% Section: introduction
%READERNOTE{SOL6-A330: दोन्ही क्रमवर्ती तक्त्यांतील असत्याचे स्वतंत्र स्वयंसिद्धक प्रस्तावनेतही स्पष्ट केले; बहुसंचाचे सर्व अर्थ जपले.}

\documentclass[../../../include/open-logic-section]{subfiles}

\begin{document}

\olfileid{pt}{seq}{int}

\olsection{परिचय}

क्रमवर्ती कलन या !!{proof} पद्धती गेरहार्ड
गेन्ट्झेनने १९३० च्या दशकात मांडल्या. व्यावहारिक
!!{proving} पद्धती म्हणून त्यांचा उद्देश तुलनेने कमी
होता. त्याऐवजी, !!{proof}s ची संभाव्य रचना आणि
त्यांचे गुणधर्म यांविषयी काही निष्कर्ष सिद्ध करण्यासाठी
गेन्ट्झेनने त्यांचा वापर केला. नेमके असेच प्रश्न
सिद्धता-उपपत्तीच्या अभ्यासाचे विषय आहेत.

नावाप्रमाणेच क्रमवर्ती कलने !!{formula}s वर
नव्हे तर \emph{क्रमवर्तींवर} काम करतात. क्रमवर्ती
म्हणजे !!{formula}s च्या क्रमिकांची किंवा बहुसंचांची
जोडी. गेन्ट्झेनच्या मूळ पद्धतींमध्ये (\Log{LK} आणि
\Log{LJ}) क्रमिका वापरल्या होत्या; आता
सिद्धता-उपपत्तीचे अभ्यासक बहुसंच वापरणे पसंत करतात.
बहुसंच हा संचासारखा असतो, पण त्यात एखादा
!!{element} एकाहून अधिक वेळा येऊ शकतो.
संचाप्रमाणेच, मात्र, !!{element}s चा क्रम महत्त्वाचा
नसतो. म्हणून $\{!A, !A, !B\}$ आणि
$\{!A, !B, !A\}$ ही एकाच बहुसंचाची रूपे आहेत.
पण हा बहुसंच $\{!A,
!B\}$ आणि $\{!A, !A, !B, !B\}$ यांहून वेगळा आहे,
कारण त्यांतील $!A$ आणि $!B$ यांच्या प्रतींच्या
संख्या वेगवेगळ्या आहेत.

सिद्धता-उपपत्तीमध्ये !!{formula}s चे बहुसंच
परंपरेने मोठ्या ग्रीक अक्षरांनी लिहितात. म्हणून
$\Gamma$, $\Delta$, $\Pi$, $\Lambda$, किंवा
$\Theta$ लिहिल्यास आपण ज्या तर्कशास्त्रात काम
करतो त्यातील !!{formula}s चे बहुसंच अभिप्रेत असतात.
तसेच परंपरेने, $\tuple{\Gamma, \Delta}$ असे लिहिण्याऐवजी
सिद्धता-उपपत्तीचे अभ्यासक दोन घटक क्रमवर्ती-चिन्हाने
वेगळे करतात; येथे आपण~$\Sequent$ वापरतो.

\begin{defn}[क्रमवर्ती]
\emph{क्रमवर्ती} म्हणजे पुढील रूपाची अभिव्यक्ती:
\[
\Gamma \Sequent \Delta
\]
येथे $\Gamma$ आणि $\Delta$ हे !!{formula}s चे
सान्त (कदाचित रिकामे) बहुसंच असतात. $\Gamma$ ला
\emph{पूर्वांग}, तर $\Delta$ ला \emph{उत्तरांग} म्हणतात.
\end{defn}

$!G$ हे~$\Gamma$ मधील !!{formula}s चे संयोजन
मानू (रिकामे संयोजन~$\ltrue$); आणि $!D$ हे~$\Delta$
मधील !!{formula}s चे वियोजन मानू (रिकामे वियोजन
~$\lfalse$). मग क्रमवर्ती~$\Gamma \Sequent \Delta$
काही रचना~$\Struct{M}$ मध्ये सत्य असेल
तेव्हाच $!G \lif !D$ सत्य असते. अभिजात
तर्कशास्त्रात याचा अर्थ असा:~$\Gamma$ मधील
!!{formula}s पैकी किमान एक असत्य असते किंवा~$\Delta$
मधील !!{formula}s पैकी किमान एक सत्य असते.

$\Gamma$ हा !!{formula}s चा बहुसंच असेल
तर~$\Gamma$ मध्ये~$!A$ जोडल्यावर मिळणारा बहुसंच $\Gamma, !A$
(किंवा $!A, \Gamma$) असा लिहितो. $\Delta$ हाही
!!{formula}s चा बहुसंच असेल तर $\Gamma, \Delta$
म्हणजे त्या दोन बहुसंचांचे बहुसंच-संयोजन: $\Gamma$
मध्ये $!A$ ची \emph{प्रतींची संख्या}~n (म्हणजे~$!A$
च्या $n$ प्रती) आणि $\Delta$ मध्ये ती~$m$ असेल,
तर $\Gamma, \Delta$ मध्ये $!A$ ची प्रतींची संख्या
~$n+m$ असते. क्रमवर्तीच्या एका बाजूचा बहुसंच रिकामा
असेल तर बहुतेक वेळा ती जागा कोरी ठेवतो. उदाहरणार्थ,
$\Sequent !A$ या क्रमवर्तीचे पूर्वांग रिकामे असून
उत्तरांगात~$!A$ ची नेमकी~$1$ प्रत असते.

क्रमवर्ती कलनातील !!^a{proof} म्हणजे
क्रमवर्तींचा वृक्ष: त्याच्या सर्वांत वरच्या (किंवा
प्रारंभीच्या) क्रमवर्ती संबंधित पद्धतीतील स्वयंसिद्धके
असतात; आणि एखादी क्रमवर्ती एक किंवा दोन क्रमवर्तींच्या
खाली असेल, तर ती पद्धतीतील निगमन नियमाने त्यांच्यापासून
योग्य रीतीने मिळाली पाहिजे. प्रथम आपण अभिजात
तर्कशास्त्राच्या दोन भिन्न क्रमवर्ती पद्धती, $\Log{G1c}$
आणि~$\Log{G3c}$, पाहू. त्यांची स्वयंसिद्धके आणि नियम
तक्ते~\ref{pt:seq:int:tab:G1c} आणि~\ref{pt:seq:int:tab:G3c}
मध्ये दिले आहेत. \oliflabeldef{fol:seq::chap}{ (गेन्ट्झेनची
मूळ पद्धत \Log{LK} चे वर्णन \olref[fol][seq][]{chap}
मध्ये आहे.)}{}

\subfile{rules-G1c}
\subfile{rules-G3c}

या पद्धतींमध्ये सूक्ष्म फरक आहेत आणि त्यामुळे
त्यांचे सिद्धता-उपपत्तीय गुणधर्म वेगवेगळे आहेत.
\Log{G1c} ची स्वयंसिद्धके $!A
\Sequent !A$ अशी असतात; त्यात इतर कोणतेही
!!{formula}s चालत नाहीत. \Log{G3c} ची स्वयंसिद्धके
$!C, \Gamma \Sequent \Delta, !C$ या रूपाची
असतात; त्यात~$\Gamma$ आणि~$\Delta$ मधील इतर
``बाजूची'' !!{formula}s मनमानी असू शकतात, पण
दोन्ही बाजूंना येणारे~$!C$ हे !!{formula}
\emph{अण्वीय} असले पाहिजे. \Log{G1c} मध्ये
$\LeftR{\land}$ आणि $\RightR{\lor}$ नियमांच्या
प्रत्येकी दोन आवृत्त्या आहेत, तर \Log{G3c} मध्ये
प्रत्येकी एकच आहे. तसेच \Log{G1c} मध्ये
``आकुंचन'' (\LeftR{\Contraction}, \RightR{\Contraction})
आणि ``दुर्बलीकरण'' (\LeftR{\Weakening}, \RightR{\Weakening})
हे अतिरिक्त ``संरचनात्मक नियम'' आहेत.
यांशिवाय असत्याचे स्वयंसिद्धकही आहे: \Log{G1c} मध्ये
$\lfalse\Sequent\quad$, आणि \Log{G3c} मध्ये
$\lfalse,\Gamma\Sequent\Delta$. त्यामुळे समान
अणुरूप सूत्र दोन्ही बाजूंना असण्याचा प्रसंग हाच
स्वयंसिद्धकांचा एकमेव प्रसंग नाही.

\Log{G1c} आणि \Log{G3c} या दोन्ही
पद्धती अभिजात तर्कशास्त्रासाठी निर्दोष व संपूर्ण
आहेत. म्हणजे या पद्धतींमध्ये !!{provable} असलेली
प्रत्येक क्रमवर्ती प्रत्येक !!{structure} मध्ये सत्य
असते; आणि प्रत्येक !!{structure} मध्ये सत्य असलेल्या
प्रत्येक क्रमवर्तीची !!a{proof} असते. म्हणून एखाद्या
क्रमवर्तीची !!a{proof} \emph{आहे का}, या प्रश्नाचे
उत्तर तर्कशास्त्राच्या इतर पद्धतींनीही (उदा.,
प्रतिमान-उपपत्तीने) मिळवता येते. आणि दोन्ही पद्धती
नेमक्या प्रत्येक !!{structure} मध्ये सत्य असलेल्या
क्रमवर्ती !!{prove} करतात; म्हणून त्या विशेषतः
त्याच क्रमवर्ती सिद्ध करतात.

पण सिद्धता-उपपत्तीला एखादी गोष्ट
\emph{सिद्धतायोग्य आहे का} एवढाच नव्हे तर ती
\emph{कशी} सिद्ध होते हाही प्रश्न महत्त्वाचा असतो.
सिद्धता-उपपत्तीचे अभ्यासक असे प्रश्न विचारतात:
``एखादा विशिष्ट नियम नेहमी टाळता येईल का?'',
``नव्या क्रमवर्ती सिद्धतायोग्य न होता हा नियम पद्धतीत
जोडता येईल का?'', किंवा ``एका पद्धतीतील !!{proof}s
दुसऱ्या पद्धतीतील !!{proof}s मध्ये रूपांतरित करता
येतील का?'' अशा प्रश्नाचे उत्तर ``होय'' असेल, तर
त्या तथ्याची सिद्धता बहुधा !!{proof}s च्या
गुंतागुंतीवर किंवा, समतुल्यपणे, त्या !!{proof}s च्या रचनेवर
विगमन करून दिली जाते. त्यातून केवळ तथ्याचीच
सिद्धता मिळत नाही (उदा., \Log{G1c} ने एखादी
क्रमवर्ती सिद्ध केली तर तीच क्रमवर्ती \Log{G3c}
नेही सिद्ध होते), तर एक प्रक्रिया मिळते (उदा.,
\Log{G1c} मधील !!{proof}s ना \Log{G3c} मधील
!!{proof}s मध्ये रूपांतरित करण्याची).

सिद्धता-उपपत्तीतील एक मध्यवर्ती निष्कर्ष म्हणजे
\emph{कट-निरसन प्रमेय}. त्यानुसार पुढील कट नियम
\[
\Axiom$\Gamma \fCenter \Delta, !A$
\Axiom$!A, \Pi \fCenter \Lambda$
\RightLabel{\Cut}
\BinaryInf$\Gamma, \Pi \fCenter \Delta, \Lambda$
\DisplayProof
\]
क्रमवर्ती पद्धतींमध्ये जोडला तरी कोणत्या क्रमवर्ती
!!{provable} आहेत यांत बदल होत नाही. याहूनही,
त्या निष्कर्षाची सिद्धता अशी पद्धत देते की
\Log{G1c} (किंवा \Log{G3c}) चे नियम आणि~\Cut{}
वापरून केलेल्या कोणत्याही !!{proof} चे रूपांतर
फक्त \Log{G1c} (किंवा \Log{G3c}) चेच नियम
वापरणाऱ्या सिद्धतेत करता येते. दुसऱ्या शब्दांत, ही
प्रक्रिया~\Cut{} वापरणाऱ्या !!{proof}s मधून
हा नियम ``निरसित'' करते; म्हणून हे नाव.

\end{document}
